On étudie le fonctionnement d'un système de chauffage installé dans une pièce.

Ce système se déclenche automatiquement dès que la température de la pièce est inférieure ou égale à 18\,\textdegree{}C (degrés Celsius), et s'éteint lorsqu'elle atteint 20\,\textdegree{}C.

	\subsection*{Partie A : Phase de chauffage}
		Pour une température de la pièce variant de 18\,\textdegree{}C à 20\,\textdegree{}C, le système de chauffage fonctionne en continu. La température de la pièce augmente progressivement.

		Dans cette partie, on modélise la température de la pièce, en degré Celsius, à l'instant $t$, exprimé en dizaines de minutes, par une fonction $T$ définie sur $[0 ;+\infty[$.

		On admet que la fonction $T$ est :
		\begin{itemize}[label={\textbullet~}]
			\item dérivable sur $[0 ;+\infty[$;
			\item solution de l'équation différentielle
			\[(E): \quad y'=-0,035 y+0,91\]
			où $y$ est une fonction de la variable $t$, définie et dérivable sur $[0 ;+\infty[$.
		\end{itemize}

		On note $T'$ la fonction dérivée de la fonction $T$.

		On suppose qu'au début de l'étude, la température de la pièce est de 18\,\textdegree{}C.

		Ainsi \quad $T(0)=18$.

		\begin{enumerate}
			\item Donner les solutions de l'équation différentielle $(E)$ sur $[0 ; +\infty[$.

			\item En déduire que, pour tout nombre réel $t$ appartenant à l'intervalle $[0 ;+\infty[$, on a :
			\[ T(t)=26-8 \e^{-0,035 t}.\]

			\item Selon ce modèle, déterminer au bout de combien de temps la pièce atteindra la température de 20\,\textdegree{}C.

			On exprimera le résultat en heures et minutes arrondi à la minute.

			\item Si une panne du système de chauffage l'empêche de s'éteindre lorsque la température de la pièce atteint 20\,\textdegree{}C, la température de la pièce pourra-t-elle dépasser 28\,\textdegree{}C selon ce modèle ?

			Justifier.
		\end{enumerate}

	\subsection*{Partie B: Phase de refroidissement}
		Lorsque la pièce atteint la température de 20\,\textdegree{}C, le système de chauffage s'éteint et la pièce refroidit.

		On modélise la température de la pièce par la suite $(u_{n})$ définie par $u_{0}=20$ et pour tout entier naturel~$n$ :
		\[u_{n+1}=0,965 u_{n}+0,35+0,07 \e^{-0,1 n}\,,\]
		où $n$ est un nombre entier exprimant le temps écoulé en dizaines de minutes.

		\begin{enumerate}
			\item Montrer que \quad $u_{1}=19,72$.

			\item Montrer par récurrence que, pour tout entier naturel $n$, on a
			\[u_{n}>10.\]
		\end{enumerate}

		On admet que la suite $(u_{n})$ est décroissante.

		\begin{enumerate}[resume]
			\item En déduire que la suite $(u_{n})$ est convergente.

			\item On note $\ell$ la limite de la suite $\left(u_{n}\right)$.

			\begin{enumerate}
				\item Justifier que cette limite est solution de l'équation
				\[x=0,965 x+0,35.\]

				\item Déterminer $\ell$ et interpréter le résultat dans le contexte de l'exercice.
			\end{enumerate}

			\item On rappelle que le système de chauffage se met en marche automatiquement dès que la température de la pièce, est inférieure ou égale à 18\,\textdegree{}C.
			\begin{enumerate}
				\item Recopier et compléter les lignes 4,5 et 6 du programme écrit en langage Python ci-dessous, afin qu'il renvoie le nombre de dizaines de minutes à partir duquel le système de chauffage se remettra en marche.
				\begin{center}
					\begin{tblr}{vline{1,3}={solid}, hline{1,Z}={solid},colspec={X[r,3mm] X[4cm,font=\ttfamily]},rowsep=0pt,column{1}={bg=gray9}}
						1.& def marche() :\\
						2.& \quad n = 0\\
						3.& \quad u = 20\\
						4.& \quad while \dots :\\
						5.& \quad \quad u = \dots\\
						6.& \quad \quad n = \dots\\
						7.& \quad return n\\
					\end{tblr}
				\end{center}

				\item Déterminer le temps, en dizaines de minutes, à partir duquel le système de chauffage se remettra en marche.
			\end{enumerate}
		\end{enumerate}
